% Part: sets-functions-relations
% Chapter: arithmetization
% Section: reflections

\documentclass[../../../include/open-logic-section]{subfiles}

\begin{document}

\olfileid[pt]{sfr}{arith}{ref}
\olsection{Algumas Reflexões Filosóficas}

Quanto aos pormenores técnicos, por ora basta. Mas o que foi que alcançaram?

Bem, de forma pouco contestável, deram-nos uma matemática pura \emph{belíssima}.
Além disso, houve conquistas conceptuais profundas. Foi um insight
profundo perceber que a propriedade de completude exprime a diferença crucial
entre os reais e os racionais. Além disso, a construção explícita dos reais,
como cortes de Dedekind, põe o objeto de estudo da análise em bases sólidas.
Sabemos que a noção de \emph{corpo ordenado completo} é coerente, pois os cortes
formam precisamente um tal corpo.

Todavia, importa também levantar algumas reservas quanto a essas conquistas.

Primeiro, não é claro que pensar os reais em termos de cortes seja
\emph{mais} rigoroso do que pensá-los nas suas expansões decimais familiares
(eventualmente infinitas). Esta última ``construção'' dos reais tem alguma
semelhança com a construção via sequências de Cauchy; mas, de fato, era,
essencialmente, conhecida dos matemáticos desde princípios do século XVII (ver
\olref[cauchy]{sec}). O verdadeiro aumento de rigor veio da tomada de consciência
de que os reais têm a propriedade de completude; a possibilidade de construir
números reais como conjuntos particulares talvez não seja, por si só, tão
interessante.

É ainda menos claro que a aritmetização (muito mais fácil) dos inteiros ou dos
racionais aumente o rigor nessas áreas. Aqui, vale a pena uma observação simples.
Tendo \emph{construído} os inteiros como classes de equivalência de pares
ordenados de naturais, e depois construído os racionais como classes de
equivalência de pares ordenados de inteiros, e depois construído os reais como
conjuntos de racionais, \emph{nos esquecemos de imediato} das construções. Em
particular: ninguém quereria \emph{invocar} estas construções no decurso de uma
prova matemática (exceto, claro, numa prova de que as construções se comportam
como deviam). É muito mais fácil falar de um real diretamente do que falar de
algum conjunto de conjuntos de conjuntos de conjuntos de conjuntos de conjuntos
de conjuntos de naturais. 
%(For reals were constructed as sets of rationals; and these were constructed as equivalence classes of ordered pairs of integers, O.s.\@, sets of sets of sets of integers; etc.).
%2/3 = [2,3]
%	O.s.\@, {<2,3>, ... }
%	{ {{2},{2,3}}, ... }
%
%reals are 
%	sets of rationals 
%
%rationals are 
%	sets of sets of sets of integers
%
%integers are
%	sets of sets of sets of naturals
	
O mais duvidoso, porém, é que estas definições nos digam o que os inteiros, racionais
ou reais \emph{são}, \emph{do ponto de vista metafísico}. Isto é, é duvidoso que
os reais (digamos) \emph{sejam} certos conjuntos (de conjuntos de
conjuntos\ldots). A principal barreira a tal visão é que a construção poderia ter
sido feita de muitas maneiras diferentes. No caso dos reais, há construções
genuinamente diferentes e interessantes (ver \olref[cauchy]{sec}). Mas eis uma
forma trivial de obter construções diferentes: como em
\olref[sfr][rel][ref]{sec}, poderíamos ter definido pares ordenados de modo
ligeiramente diferente; se tivéssemos usado essa noção alternativa de par
ordenado, as nossas construções teriam funcionado exatamente tão bem como
funcionaram, mas teríamos acabado com objetos diferentes. Assim, há muitas
construções rivais, em teoria dos conjuntos, dos inteiros, racionais e reais. Seria,
ademais, arbitrário (e embaraçoso) afirmar que os inteiros (digamos) são
\emph{estes} conjuntos, e não \emph{aqueles}. (Como na \olref[sfr][rel][ref]{sec},
isto é um caso de um argumento tornado célebre por \citealt{Benacerraf1965}.)

Um outro ponto merece ser levantado: há algo de bastante \emph{estranho} nas
nossas construções. Começamos pelos números naturais. Depois, construímos os
inteiros e construímos ``o $0$ dos inteiros'', O.s.\@, $\equivrep{0,0}{\Intequiv}$.
Mas $0 \neq \equivrep{0,0}{\Intequiv}$. De fato, dadas as nossas construções,
\emph{nenhum} natural é um inteiro. Isto parece fortemente contraintuitivo.
Na \olref[sfr][set][imp]{sec}, afirmamos, sem grande argumento, que
$\Nat \subseteq \Rat$. Agora, se as construções nos dissessem exatamente \emph{o que}
são os números, esta afirmação seria trivialmente falsa. 

Recuando um pouco, onde chegamos? Trabalhando numa teoria ingênua dos conjuntos e
tomando os naturais como dados, podemos \emph{tratar} inteiros, racionais e reais
como certos conjuntos. Nesse sentido, podemos \emph{embutir} as teorias destas
entidades numa teoria dos conjuntos. Mas a importância filosófica deste embutimento
não é assim tão direta.

Claro que isto não é a última palavra!{} O ponto é apenas este: mostrar que a
aritmetização dos reais \emph{é de} profunda significação filosófica exigiria
argumento filosófico adicional.

%Additionally: for the entire duration of this chapter, we have treated
%the natural numbers as simply \emph{given} to us. But what can we do
%with them? That is the subject matter for the next chapter.

\end{document}
